\chapter{附录 C：Matlab 代码注释示例}

\section{引言}

本附录给出一个带注释的 Matlab 代码示例，展示如何借助 SPM12 中提供的标准反演与绘图例程来指定并求解一个生成模型。这里复现的是第 7 章中的 T 迷宫觅食示例。单独阅读这一附录可能会略显枯燥，但如果读者尝试在 Matlab 中亲自实现并运行这段代码，它会非常有帮助，因为只有这样，主动推断的机械过程才会真正变得直观。我们建议读者通过修改不同参数值、甚至有意“破坏”这个 demo 来加深理解。

\section{预备}

我们假定读者对 Matlab 已有基本熟悉，并已经成功从 \url{https://www.fil.ion.ucl.ac.uk/spm/} 下载了 SPM12 软件包。第一步，是确保包含 SPM12 函数的文件夹已经加入 Matlab 的搜索路径。随后，我们打开一个 Matlab 脚本，先定义一个函数并为其命名（这里命名为 \verb|demo_AI_book|）：

\begin{verbatim}
function demo_AI_book
rng default
\end{verbatim}

\noindent\textbf{图 C.1}

图 C.1 中第二行的作用，是把随机数生成器设置为默认初始种子，从而保证每次调用该函数时生成同一组随机数。我们通常会在 demo 中这样做，以确保结果可复现。不过，如果我们想多次运行该函数，以计算跨多次试次行为的汇总统计量，也可以省略这一步。一般来说，最好在这里附上一些说明脚本用途的注释；但由于本附录本身就是对该脚本的逐段注释，这里便不再重复。

接下来，我们定义稍后会用到的一些重要常数。把这些量集中写在一起的好处是，后面如果想扰动它们、测试模型敏感性，就能很容易找到。这里我们定义两个将扮演概率角色的参数，它们在 C.3 节中会变得清晰起来：

\begin{verbatim}
a = .98;
b = 1 - a;
\end{verbatim}

\noindent\textbf{图 C.2}

这样的定义保证了 $a+b=1$。此时，我们就可以开始设置构成生成模型的 $A$、$B$、$C$、$D$ 矩阵与向量。在这个简单仿真中，我们假定这些参数在生成模型与生成过程里是相同的。

\section{似然}

这里我们的重点，是说明如何在 Matlab 里形式化一个似然矩阵，因此不会花太多篇幅去重复描述该生成模型及其对应范式本身；相关背景可参见第 7.2 节。我们的目标，是把图 7.4 与图 7.5 中给出的似然矩阵，翻译成反演例程能够理解的语法。我们从 $A_1$ 开始，它在 Matlab 中写作 \verb|A{1}|，其中花括号里的项对应上标，也就是结果模态。矩阵（更严格地说是张量）的元素可以用三个索引访问，分别对应结果、第一个隐藏状态因子（位置）以及第二个隐藏状态因子（上下文）。吸引性刺激位于右侧的上下文，可在第三个位置上用索引 $1$ 指定；而当其位于左侧时，则用该位置上的索引 $2$ 指定：

\begin{verbatim}
A{1} (:, :, 1) = [...
1 0 0 0; % start
0 0 0 0; % left cue
0 1 0 0; % right cue
0 0 1 0 % left
0 0 0 1]; % right
A{1} (:, :, 2) = [...
1 0 0 0; % start
0 1 0 0; % left cue
0 0 0 0; % right cue
0 0 1 0 % left
0 0 0 1]; % right
\end{verbatim}

图 C.3

行表示结果，列表示第一个隐藏状态因子的不同可能水平。将其与图 7.4 和图 7.5 中展示的矩阵对照，有助于理解这种语法。$A_2$ 矩阵也以类似方式为上下文加以定义。这里，我们使用脚本顶部在图 C.4 中定义的 $a$ 和 $b$：

\begin{verbatim}
A{2} (:, :, 1) = [...
1 1 0 0; % reward neutral
0 0 a b; % reward positive
0 0 b a]; % reward negative
A{2} (:, :, 2) = [...
1 1 0 0; % reward neutral
0 0 b a; % reward positive
0 0 a b]; % reward negative
\end{verbatim}

图 C.4

至此，我们完成了对 $A$ 的设定。对于两个模态中的每一种结果组合（以上标或花括号区分，对应各行），以及两个隐藏状态因子的每一种取值组合（由第二和第三个索引区分），我们都已经定义了相应的概率。

\subsection{转移概率}

在 $A$ 之后，我们接着指定 $B$。如第 7 章所讨论，与 $B$ 矩阵相关联的上标指的是隐藏状态因子，而不是像 $A$ 那样表示结果模态；同样，我们仍将花括号视为上标的等价写法（图 C.5）。回忆一下，这两个因子分别是位置（1）和上下文（2）。每个矩阵都把前一时刻的状态（列）映射到当前时刻的状态（行）。如果某个状态因子是可控制的，那么 $B$ 矩阵还可以随动作而变化。这意味着控制状态需要额外一个索引，用来指明采取的是哪一个动作：

\begin{verbatim}
B{1}(:, :, 1)     = [1 1 0 0; 0 0 0 0;0 0 1 0;0 0 0 1];
B{1}(:, :, 2)     = [0 0 0 0; 1 1 0 0;0 0 1 0;0 0 0 1];
B{1}(:, :, 3)     = [0 0 0 0; 0 0 0 0;1 1 1 0;0 0 0 1];
B{1}(:, :, 4)     = [0 0 0 0; 0 0 0 0;0 0 1 0;1 1 0 1];

B{2} =
eye(2);
\end{verbatim}

图 C.5

这里，我们针对四个可用动作中的每一个，给出了可控制的位置状态。回忆一下，一共有四个可用动作，每个动作都决定向四个位置之一发生转移。右臂和左臂都是吸收态，这意味着老鼠一旦进入这些位置，就必须停留在那里。这里的分号表示矩阵中每一行的结束。上下文状态并不受生物体控制，因此我们不需要为每个动作分别指定一组不同的转移概率。我们只需定义一个单一的单位矩阵，在 Matlab 中记作 \verb|eye|。这就实现了图 7.6 和公式 7.5 中的设定。

\subsection{先验偏好与初始状态}

在 $B$ 之后是 $C$。这里我们回到了与 $A$ 相似的语法：上标和花括号都对应结果模态。此外，我们还像处理 $B$ 那样指定 $D$：上标和花括号对应的是隐藏状态因子（图 C.6）。$C$ 和 $D$ 的行数必须分别与对应的 $A$ 矩阵和 $B$ 矩阵的行数一致。

\begin{verbatim}
C{1}   = [-1 -1 -1;
0 0 0;
0 0 0;
0 0 0;
0 0 0];
c      = 6;
C{2}   = [ 0 0 0;
c c c;
-c -c -c];
D{1} = [1 0 0 0]';
D{2} = [1 1]'/2;
\end{verbatim}

图 C.6

偏好是在 $C$ 矩阵中以对数概率的形式给出的（它们不必归一化）。两个结果之间若相差 6，就意味着概率更高的那个结果比另一个结果高出 $\exp(6)$ 倍。每一行对应一个不同的结果，列则对应不同的时间步。这使得时间变化的偏好成为可能。如果只给出一列，就默认每个时间点上的偏好都相同。这里给出的 $D$ 向量，则只是为试次开始时的每个状态赋予一个概率。可将其与公式 7.6 和公式 7.7 对照理解。还要注意，在 Matlab 中，转置记号写作 \verb|'|。

\subsection{策略空间}

我们已经通过 $B$ 矩阵隐式地指定了允许的动作。我们可以假设策略与动作完全等同，并且在每个时间步都重新选择策略。另一种做法是把策略指定为动作序列。可以这样理解：每个策略都告诉我们，在每个时间步上，究竟是哪一个 $B$ 矩阵（由第三个索引标识）在起作用。我们通过指定一个数组 $V$ 来实现这一点（图 C.7）：

\begin{verbatim}
V(:, :, 1) = [1 1 1 1 3 4 2 2 2 2
1 2 3 4 3 4 1 2 3 4];
V(:, :, 2) = 1;
\end{verbatim}

图 C.7

$V$ 的第一个索引（即各行）表示动作序列中的位置。这意味着第一行是第一个动作，第二行是第二个动作。由于动作会导致状态转移，一个三步模型只会

\section{整合各个组成部分}

在完成 POMDP 的各项设定之后，我们现在将其整合到一个单独的 \verb|mdp| 变量中（见图 C.8）：

\begin{verbatim}
mdp.V = V; % allowable policies
mdp.A = A; % observation model
mdp.B = B; % transition probabilities
mdp.C = C; % preferred outcomes
mdp.D = D; % prior over initial states
mdp.S = [1 1]'; % true initial state
\end{verbatim}

\noindent\textbf{图 C.8}

这里最后一行使我们能够指定希望作为起点的真实隐藏状态。这个信息对我们模拟出来的生物体而言是不可用的；它必须依据这些状态所生成的结果，以及它自身的生成模型（由前五行给出），来推断状态。

\section{仿真与绘图}

\verb|spm_MDP_VB_X| 函数在幕后完成了主要工作；它实现了第 4 章和第 7 章中描述的消息传递与策略选择。此外，它还会模拟我们的生物体必须面对的世界，包括状态之间的转移以及结果的生成。我们本来还可以加入许多其他选项；有关细节，请读者参阅 Matlab 脚本中该函数的说明文档。

一旦完成一次试次的仿真，我们通常会希望找到某种图形化方式来呈现结果。像图 7.2 和图 7.7 那样的绘图，可以通过标准绘图例程自动生成：

\begin{verbatim}
MDP = spm_MDP_VB_X(mdp);
spm_figure('GetWin', 'Figure 1'); clf
spm_MDP_VB_trial(MDP);

spm_figure('GetWin', 'Figure 2'); clf
spm_MDP_VB_LFP(MDP,[],1);

spm_figure('GetWin','Figure 3'); clf
spm_MDP_VB_LFP(MDP,[],2);
\end{verbatim}

\noindent\textbf{图 C.9}

图 C.9 中的这些代码行会完成仿真，并将结果绘制在三个图中。第一幅图给出仿真的总体摘要，包括状态、结果、被选中的策略，以及回顾性推断。第二幅图展示第一个隐藏状态因子在信念更新过程中的电生理相关量；第三幅图则对第二个因子做同样的展示。为了让你确信一切都在正常运行，我们把前几行代码的输出以图形形式重现出来（见图 C.10）。我们本来也可以把另外两幅图也重现出来；不过，我们希望给你一个亲自运行脚本的机会，通过实践主动推断来降低自己对这些图究竟展示了什么的不确定性。

这个带注释的简单示例到这里就结束了。我们希望它能引导你开始探索那些可以用相同原则和语法来指定的各种生成模型。你可以在 Matlab 命令行中输入 \verb|DEM|，然后在随后出现的图形用户界面中选择相应的演示，以找到更多示例。我们提供的 demo 脚本，是 Friston、FitzGerald 等人（2017）所使用例程的一个简化版本。若想进一步了解更多细节，包括如何在多次试次中学习生成模型，我们建议读者参阅该论文，以及 SPM12 中提供的 \verb|DEM_demo_MDP_X.m| 脚本。

\noindent\textbf{图 C.10}\quad \verb|spm_MDP_VB_trial| 例程生成的输出。左上：对每个隐藏状态因子的回顾性信念（即在所有观测都已经获得之后）。黑色阴影表示概率为 1，白色表示概率为 0。灰色圆点（在彩色绘图中会显示为青色）表示由模拟环境生成的真实状态；我们可以看到，模拟出来的大鼠能够准确且有信心地推断出它的位置（因子 1）以及情境（因子 2）。左中：C.6 节中 \verb|V| 变量所指定的“可选策略”，每一行表示一个策略：第一列是第一步采取的动作，第二列是第二步采取的动作。每个元素中的不同阴影代表可被选择的不同动作。左下：每个模态中的真实结果（灰色圆点）。背景阴影显示的是每个模态的 \verb|C| 矩阵，颜色越深表示该结果越受偏好。右上：被推断并被选择的动作。右中：每个时间步对各个策略的信念。右下：关于策略信念的推断精度（灰线），以及其变化率的柱状图；这种形式让人联想到电生理学中用于展示多巴胺能神经元放电的栅格图。注意，这里的时间是以更新次数来表示的；这些更新对应于对自由能进行梯度下降时的迭代。默认情况下，每个时间步有 16 次更新。
